% part: intuitionistic-logic
% chapter: introduction
% section: axiomatic-derivations

\documentclass[../../../include/open-logic-section]{subfiles}

\begin{document}

\olfileid{int}{int}{axd}

\olsection{স্বতঃসিদ্ধিক \usetoken{P}{derivation}}

স্বজ্ঞাবাদী বচনমূলক যুক্তিবিদ্যার স্বতঃসিদ্ধিক !!{derivation}s ধারণাগতভাবে
সবচেয়ে সরল এবং ঐতিহাসিকভাবে প্রথম !!{derivation} ব্যবস্থা। এগুলি
ধ্রুপদি বচনমূলক যুক্তিবিদ্যার মতোই কাজ করে।

\begin{defn}[!!^{derivability}]
$\Gamma$ যদি $\Lang L$-এর !!{formula}s-এর একটি সেট হয়, তবে
$\Gamma$ থেকে একটি \emph{!!{derivation}} হলো !!{formula}s-এর একটি সসীম
ক্রম $!A_1$, \dots,~$!A_n$, যেখানে প্রতিটি $i \le n$-এর জন্য নিচের
কোনো একটি শর্ত পূরণ হয়:
\begin{enumerate}
\item $!A_i \in \Gamma$; অথবা
\item $!A_i$ একটি স্বতঃসিদ্ধ; অথবা
\item $j < i$ ও $k < i$ এমন কোনো $!A_j$ এবং $!A_k$ থেকে মোডাস পোনেন্স
  নিয়মে $!A_i$ পাওয়া যায়; অর্থাৎ, $!A_k \ident !A_j \lif !A_i$।
\end{enumerate}
\end{defn}

\begin{defn}[স্বতঃসিদ্ধ]
স্বজ্ঞাবাদী বচনমূলক যুক্তিবিদ্যার \emph{স্বতঃসিদ্ধ}গুলির সেট~$\PAx$-এ
নিচের রূপগুলির সব !!{formula}s আছে:
\begin{align}
  & (!A \land !B) \lif !A \ollabel{ax:land1}\\
  & (!A \land !B) \lif !B \ollabel{ax:land2}\\
  & !A \lif (!B \lif (!A \land !B)) \ollabel{ax:land3}\\
  & !A \lif (!A \lor !B) \ollabel{ax:lor1}\\
  & !A \lif (!B \lor !A) \ollabel{ax:lor2}\\
  & (!A \lif !C) \lif ((!B \lif !C) \lif ((!A \lor !B) \lif !C)) \ollabel{ax:lor3}\\
  & !A \lif (!B \lif !A) \ollabel{ax:lif1}\\
  & (!A \lif (!B \lif !C)) \lif ((!A \lif !B) \lif (!A \lif !C)) \ollabel{ax:lif2}\\
  & \lfalse \lif !A \ollabel{ax:lfalse1}
\end{align}
\end{defn}

\begin{defn}[!!^{derivability}]
!!^a{formula}~$!A$-কে $\Gamma$ থেকে \emph{!!{derivable}} বলা হয়,
লিখি $\Gamma \Proves !A$, যদি $\Gamma$ থেকে এমন একটি !!a{derivation}
থাকে যার শেষ সূত্র $!A$।
\end{defn}

\begin{defn}[উপপাদ্য]
!!^a{formula}~$!A$ একটি \emph{উপপাদ্য}, যদি শূন্য সেট থেকে~$!A$-এর
একটি !!a{derivation} থাকে। $!A$ উপপাদ্য হলে লিখি $\Proves !A$,
আর উপপাদ্য না হলে লিখি $\Proves/ !A$।
\end{defn}

\begin{prop}
  স্বজ্ঞাবাদী যুক্তিবিদ্যার স্বতঃসিদ্ধিক !!{derivation} ব্যবস্থায়
  $\Gamma \Proves !A$ হলে ধ্রুপদি যুক্তিবিদ্যাতেও
  $\Gamma \Proves !A$। বিশেষত, $!A$ স্বজ্ঞাবাদী উপপাদ্য হলে সেটি
  ধ্রুপদি উপপাদ্যও।
\end{prop}

\begin{proof}
  স্বজ্ঞাবাদী প্রতিটি স্বতঃসিদ্ধ ধ্রুপদি স্বতঃসিদ্ধও। তাই
  স্বজ্ঞাবাদী যুক্তিবিদ্যার প্রতিটি !!{derivation} ধ্রুপদি
  যুক্তিবিদ্যাতেও !!a{derivation}।
\end{proof}

\end{document}
